\chapter{Embeddings and Vector Databases: Geometry, Indexing, and Internals}
\label{ch:embeddings_vector_db}

\begin{chaptergoalsbox}
\textbf{What You Will Learn:}
\begin{itemize}[leftmargin=15pt]
  \item High-dimensional metric space geometry: $L_2$ Euclidean, Cosine similarity, inner product, and the concentration of measure.
  \item The curse of dimensionality: Why distance distributions collapse in $\mathbb{R}^{1536}$ and how hyperspheres behave.
  \item Inverted File Index (IVF): Voronoi cell partitioning, k-means quantization, and $nlist/nprobe$ tuning.
  \item Hierarchical Navigable Small World (HNSW) graphs: Multi-layer skip-list graph topologies, heuristic edge selection, and greedy routing.
  \item DiskANN and NVMe-native graph traversal: Product Quantization (PQ) and Vamana graphs for billion-scale collections.
  \item Deep engine internals: Qdrant's Rust segment architecture vs pgvector's PostgreSQL page buffer integration.
\end{itemize}

\textbf{What You Should Already Understand:}
\begin{itemize}[leftmargin=15pt]
  \item Linear algebra: Vectors, dot products, matrix multiplication, and vector norms ($L_1, L_2$).
  \item Basic graph theory: Directed graphs, adjacency lists, Dijkstra / greedy search.
  \item Relational storage fundamentals: Buffer pools, write-ahead logs (WAL), and disk paging.
\end{itemize}

\textbf{Why This Topic Exists:}
Exact nearest neighbor search (flat scan) scales as $O(N \cdot D)$. For 10 million 1536-dimensional vectors, every single query requires 61.4 billion floating-point operations, taking hundreds of milliseconds on high-end hardware. Approximate Nearest Neighbor (ANN) indexing structures reduce search latency to sub-millisecond speeds at the expense of a negligible recall trade-off.

\textbf{Motivating Practical Problem:}
A cybersecurity threat-intelligence platform receives 50 million streaming endpoint telemetry vectors daily. Security analysts require sub-15ms vector similarity queries against historical telemetry while simultaneously filtering by \texttt{tenant\_id}, \texttt{operating\_system}, and \texttt{timestamp \textgreater{}= NOW() - INTERVAL '7 DAYS'}. Naive vector databases crash under memory exhaustion or suffer severe post-filtering scan degradation.
\end{chaptergoalsbox}

\section{High-Dimensional Vector Geometry}

Dense embeddings project unstructured semantic entities (words, sentences, code, images) into a continuous vector space $\mathbb{R}^D$, typically with dimensionality $D \in [256, 3072]$.

\subsection{Metric Distances and Similarities}
Let $\vec{u}, \vec{v} \in \mathbb{R}^D$ be two vectors. The three principal similarity functions are:

\begin{enumerate}[leftmargin=18pt]
  \item \textbf{Euclidean Distance ($L_2$ Norm)}:
  \begin{equation}
  d_{L_2}(\vec{u}, \vec{v}) = \|\vec{u} - \vec{v}\|_2 = \sqrt{\sum_{i=1}^D (u_i - v_i)^2}
  \end{equation}
  Represents direct geometric distance. Sensitive to vector magnitude.

  \item \textbf{Cosine Distance}:
  \begin{equation}
  S_{\cos}(\vec{u}, \vec{v}) = \frac{\vec{u} \cdot \vec{v}}{\|\vec{u}\|_2 \|\vec{v}\|_2} = \frac{\sum_{i=1}^D u_i v_i}{\sqrt{\sum u_i^2} \sqrt{\sum v_i^2}}
  \end{equation}
  \begin{equation}
  d_{\cos}(\vec{u}, \vec{v}) = 1 - S_{\cos}(\vec{u}, \vec{v})
  \end{equation}
  Measures the angular deviation $\theta$ between vectors, completely invariant to length.

  \item \textbf{Inner Product (Dot Product)}:
  \begin{equation}
  \langle \vec{u}, \vec{v} \rangle = \vec{u} \cdot \vec{v} = \sum_{i=1}^D u_i v_i
  \end{equation}
\end{enumerate}

\begin{intuitionbox}[The Normalization Identity]
If all vectors are pre-normalized to unit length ($\|\vec{u}\|_2 = 1$), Euclidean distance, Cosine distance, and Inner Product become mathematically equivalent:
\begin{equation}
\|\vec{u} - \vec{v}\|_2^2 = \|\vec{u}\|^2 + \|\vec{v}\|^2 - 2(\vec{u} \cdot \vec{v}) = 1 + 1 - 2\cos(\theta) = 2(1 - \cos(\theta))
\end{equation}
Consequently:
\begin{equation}
d_{L_2}^2 = 2 \cdot d_{\cos} = 2(1 - \langle \vec{u}, \vec{v} \rangle)
\end{equation}
In production vector databases, \textbf{always pre-normalize vectors to unit length} before ingestion. Computing a pure inner product requires only $D$ multiply-accumulate (MAC) SIMD instructions, eliminating expensive square roots and division operations.
\end{intuitionbox}

\subsection{The Curse of Dimensionality: Concentration of Measure}
Human spatial intuition is formed in $\mathbb{R}^2$ and $\mathbb{R}^3$. In $\mathbb{R}^{1536}$, geometric properties become profoundly counter-intuitive.

\begin{figure}[htbp]
\centering
\begin{tikzpicture}
\draw[->, thick, color=headingblue] (0,0) -- (9,0) node[right, font=\sffamily\footnotesize] {Distance to Neighbors};
\draw[->, thick, color=headingblue] (0,0) -- (0,3.5) node[above, font=\sffamily\footnotesize] {Probability Density $p(d)$};

% 2D Distribution (Wide)
\draw[thick, color=linkblue, domain=0.5:8.5, samples=100] 
  plot (\x, {2.5 * exp(-0.5 * ((\x - 4.5)/1.8)^2) / (1.8 * sqrt(2 * 3.1415))}) 
  node[right, font=\tiny\sffamily, color=linkblue] {Low Dimension ($D=3$)};

% High-D Distribution (Dirac Spike)
\draw[very thick, color=warningRed, domain=3.0:6.0, samples=100] 
  plot (\x, {6.5 * exp(-0.5 * ((\x - 4.5)/0.25)^2) / (0.25 * sqrt(2 * 3.1415))}) 
  node[right, font=\tiny\sffamily, color=warningRed] {High Dimension ($D=1536$)};

\end{tikzpicture}
\caption{The Concentration of Measure: As dimensionality $D \to \infty$, the distance distribution collapses into a razor-thin spike centered at $\sqrt{2D}$.}
\label{fig:concentration_of_measure}
\end{figure}

As $D \to \infty$, independent random vectors become mutually orthogonal ($\vec{u} \cdot \vec{v} \to 0$, meaning $\cos(\theta) \to 90^\circ$). Furthermore, the ratio between the distance to the nearest neighbor $d_{\min}$ and the farthest neighbor $d_{\max}$ converges to 1:
\begin{equation}
\lim_{D \to \infty} \frac{d_{\max} - d_{\min}}{d_{\min}} = 0
\end{equation}
This mathematical reality implies that standard tree partitioning indexes (such as $k$-d trees or $R$-trees) degrade to $O(N)$ linear brute-force scans in high dimensions. To achieve sub-linear query times, we must turn to graph-based and quantization-based approximate indexing.

\section{The Inverted File Index (IVF)}

The Inverted File Index (IVF) reduces search complexity by partitioning the continuous vector space into discrete \textbf{Voronoi cells} using $k$-means clustering.

\begin{figure}[htbp]
\centering
\begin{tikzpicture}[
  centroid/.style={draw=warningRed, fill=warningRed, circle, inner sep=2pt},
  point/.style={draw=headingblue, fill=headingblue!40, circle, inner sep=1.2pt},
  cell/.style={draw=gray!50, dashed, thick}
]

% Voronoi cells sketch
\draw[cell] (0,0) -- (2.5, 3.5) -- (5, 3.0) -- (4, 0.5) -- cycle;
\draw[cell] (5, 3.0) -- (8.5, 3.5) -- (7.5, 0.5) -- (4, 0.5);
\draw[cell] (0,0) -- (4, 0.5) -- (3.5, -3.0) -- (-0.5, -2.5) -- cycle;
\draw[cell] (4, 0.5) -- (7.5, 0.5) -- (8.0, -2.8) -- (3.5, -3.0);

% Centroids
\node[centroid, label={[font=\tiny\sffamily]above:$C_1$}] (c1) at (2.2, 1.8) {};
\node[centroid, label={[font=\tiny\sffamily]above:$C_2$}] (c2) at (6.2, 1.7) {};
\node[centroid, label={[font=\tiny\sffamily]below:$C_3$}] (c3) at (1.8, -1.2) {};
\node[centroid, label={[font=\tiny\sffamily]below:$C_4$}] (c4) at (5.8, -1.1) {};

% Vectors in C1
\node[point] at (1.2, 1.2) {}; \node[point] at (2.8, 2.4) {}; \node[point] at (1.5, 2.5) {};
\node[point] at (3.1, 1.1) {};

% Vectors in C2
\node[point] at (5.5, 2.2) {}; \node[point] at (7.0, 1.5) {}; \node[point] at (7.8, 2.8) {};

% Query Vector
\node[draw=intuitionTeal, fill=intuitionTeal, star, star points=5, inner sep=2.5pt, label={[font=\tiny\sffamily\bfseries, color=intuitionTeal]right:Query $\vec{q}$}] (q) at (4.2, 1.2) {};

% Probe circle
\draw[draw=intuitionTeal, thick, dotted] (q) circle (2.2cm);

\end{tikzpicture}
\caption{IVF Voronoi Partitioning: A query vector finds nearest centroids ($nprobe=2$) and only scans vectors located within those designated cells.}
\label{fig:ivf_voronoi}
\end{figure}

\subsection{Construction and Query Mechanics}
\begin{enumerate}[leftmargin=18pt]
  \item \textbf{Index Build}: Run $k$-means on training vectors to identify $C = nlist$ cluster centroids $\{\vec{c}_1, \dots, \vec{c}_k\}$. Each vector in the dataset is assigned to its nearest centroid, creating an inverted list for each cluster.
  \item \textbf{Query Execution}: Given query vector $\vec{q}$:
  \begin{enumerate}
    \item Compute distance from $\vec{q}$ to all $nlist$ centroids ($O(nlist \cdot D)$).
    \item Select the top $nprobe$ closest centroids ($nprobe \ll nlist$).
    \item Scan only the vectors stored inside the inverted lists of those $nprobe$ clusters.
  \end{enumerate}
\end{enumerate}

\begin{lstlisting}[language=Python, caption={Minimal Inverted File Index (IVF) Implementation}]
import numpy as np

class MinimalIVFIndex:
    def __init__(self, nlist: int = 100, nprobe: int = 5):
        self.nlist = nlist
        self.nprobe = nprobe
        self.centroids = None
        self.inverted_lists = {} # centroid_idx -> list of (vector_id, vector)

    def train_and_index(self, vectors: np.ndarray):
        # 1. K-means clustering (simplified random initialization)
        N, D = vectors.shape
        centroid_indices = np.random.choice(N, self.nlist, replace=False)
        self.centroids = vectors[centroid_indices].copy()
        
        for i in range(self.nlist):
            self.inverted_lists[i] = []

        # 2. Assign vectors to nearest centroid
        for vec_id, vec in enumerate(vectors):
            # Compute distance to all centroids
            dists = np.linalg.norm(self.centroids - vec, axis=1)
            nearest_centroid = np.argmin(dists)
            self.inverted_lists[nearest_centroid].append((vec_id, vec))

    def search(self, query: np.ndarray, k: int = 5) -> list[tuple[int, float]]:
        # Step 1: Find nprobe closest centroids
        centroid_dists = np.linalg.norm(self.centroids - query, axis=1)
        probed_centroids = np.argsort(centroid_dists)[:self.nprobe]
        
        # Step 2: Exhaustively search vectors within selected lists
        candidates = []
        for c_idx in probed_centroids:
            for vec_id, vec in self.inverted_lists[c_idx]:
                dist = np.linalg.norm(vec - query)
                candidates.append((vec_id, dist))
                
        # Step 3: Sort candidates and return top-k
        candidates.sort(key=lambda x: x[1])
        return candidates[:k]
\end{lstlisting}

\begin{underthehoodbox}[The IVF Recall/Latency Trade-off]
IVF performance is governed by two parameters:
\begin{itemize}
  \item $nlist$: Typically set to $\sqrt{N}$ or $4\sqrt{N}$. Higher $nlist$ reduces the number of vectors per list but increases centroid search overhead.
  \item $nprobe$: Controls recall. If $nprobe=1$, search is fast but vectors lying near cell boundaries are missed (low recall $\approx 70\%$). Increasing $nprobe \to nlist$ approaches 100\% recall but degrades back to an $O(N)$ flat scan.
\end{itemize}
\end{underthehoodbox}

\section{Hierarchical Navigable Small World (HNSW) Graphs}

HNSW (Malkov \& Yashunin, 2018) is the gold standard for in-memory approximate nearest neighbor search. It combines \textbf{Navigable Small World} graph theory with the multi-layer hierarchy of \textbf{Skip-Lists}.

\begin{figure}[htbp]
\centering
\begin{tikzpicture}[
  node/.style={circle, draw=headingblue, fill=codebg, inner sep=2.5pt},
  activenode/.style={circle, draw=warningRed, fill=warningBg, inner sep=2.5pt, line width=1pt},
  entrynode/.style={circle, draw=intuitionTeal, fill=intuitionTeal, inner sep=3pt},
  layer/.style={draw=gray!40, fill=gray!4, rounded corners=3pt, minimum width=10.0cm, minimum height=1.6cm},
  arrow/.style={-{Stealth[scale=0.8]}, thick, color=headingblue},
  route/.style={-{Stealth[scale=0.9]}, very thick, color=warningRed}
]

% Layer 2 (Sparse)
\node[layer] (l2) at (0, 3.6) {};
\node[anchor=west, font=\sffamily\bfseries\tiny, color=gray!80] at (-4.8, 4.1) {Layer 2: Top / Long-Range Highway ($p \approx 0.05$)};
\node[entrynode, label={[font=\tiny\sffamily]above:Entry Point}] (e2) at (-3.5, 3.6) {};
\node[node] (n21) at (0.5, 3.6) {};
\node[activenode] (n22) at (3.5, 3.6) {};
\draw[route] (e2) -- (n21);
\draw[route] (n21) -- (n22);

% Layer 1 (Medium)
\node[layer] (l1) at (0, 1.8) {};
\node[anchor=west, font=\sffamily\bfseries\tiny, color=gray!80] at (-4.8, 2.3) {Layer 1: Intermediate Navigation ($p \approx 0.25$)};
\node[node] (n11) at (-3.5, 1.8) {};
\node[node] (n12) at (-1.5, 1.8) {};
\node[node] (n13) at (0.5, 1.8) {};
\node[activenode] (n14) at (2.2, 1.8) {};
\node[activenode] (n15) at (3.5, 1.8) {};
\draw[arrow] (n11) -- (n12); \draw[arrow] (n12) -- (n13);
\draw[route] (n15) -- (n14);

% Layer 0 (Dense)
\node[layer] (l0) at (0, 0) {};
\node[anchor=west, font=\sffamily\bfseries\tiny, color=gray!80] at (-4.8, 0.5) {Layer 0: Ground Truth Layer (All Vectors, $M$ Edges)};
\node[node] (n01) at (-3.5, 0) {};
\node[node] (n02) at (-2.0, 0) {};
\node[node] (n03) at (-0.5, 0) {};
\node[node] (n04) at (0.8, 0) {};
\node[activenode] (n05) at (2.2, 0) {};
\node[activenode] (n06) at (3.5, 0) {};
\node[entrynode, label={[font=\tiny\sffamily, color=intuitionTeal]below:Nearest Neighbor}] (target) at (2.9, 0) {};
\draw[arrow] (n01) -- (n02); \draw[arrow] (n02) -- (n03); \draw[arrow] (n03) -- (n04);
\draw[route] (n05) -- (target);

% Vertical drop lines
\draw[dashed, color=warningRed, thick] (n22) -- (n15);
\draw[dashed, color=warningRed, thick] (n14) -- (n05);

\end{tikzpicture}
\caption{HNSW Multi-Layer Graph Traversal: Greedy routing traverses long-distance skip-edges on upper layers before dropping to denser lower layers.}
\label{fig:hnsw_layers}
\end{figure}

\subsection{Probabilistic Layer Assignment}
When a new vector $\vec{v}$ is inserted into an HNSW graph, its maximum layer $l$ is assigned using a decaying probability distribution:
\begin{equation}
l = \lfloor -\ln(\text{uniform}(0, 1)) \cdot m_L \rfloor
\end{equation}
where $m_L = \frac{1}{\ln(M)}$. This guarantees that layer 0 contains all $N$ elements, while each successive layer $l$ contains a geometrically decreasing fraction $\frac{1}{M}$ of elements, matching the logarithmic complexity of a 1D skip-list in metric space.

\subsection{Key Hyperparameters}
\begin{itemize}[leftmargin=18pt]
  \item \textbf{$M$}: The maximum number of bidirectional edges per node (typically $16 \le M \le 64$). Higher $M$ improves recall on high-dimensional or clustered data but linearly increases RAM consumption ($2 \times M \times 4 \text{ bytes}$ per vector for pointers).
  \item \textbf{$efConstruction$}: The size of the dynamic priority queue candidate list maintained during index construction (typically $100 \le efConstruction \le 400$). Higher values increase index build time exponentially but produce higher graph connectivity.
  \item \textbf{$efSearch$}: The size of the priority queue during runtime querying. A query can dynamically dial $efSearch$ up or down per request: $efSearch=32$ for high-throughput, low-latency workloads; $efSearch=256$ for ultra-high-recall applications.
\end{itemize}

\section{DiskANN: Billion-Scale Out-of-Core Search}

While HNSW is blazing fast, it has an Achilles' heel: \textbf{RAM consumption}. Storing 1 billion 768-dimensional vectors in HNSW requires:
\begin{equation}
10^9 \times (768 \times 4 \text{ bytes} + 64 \text{ edges} \times 8 \text{ bytes}) \approx 3.5 \text{ Terabytes of RAM}
\end{equation}
DiskANN (Subramanya et al., Microsoft Research, 2019) eliminates the RAM barrier by storing vectors and graph edges directly on NVMe SSDs.

\begin{figure}[htbp]
\centering
\begin{tikzpicture}[
  ram/.style={draw=intuitionTeal, fill=intuitionBg, rounded corners=2pt, font=\sffamily\scriptsize, minimum width=9.5cm, minimum height=1.4cm, align=left},
  ssd/.style={draw=headingblue, fill=codebg, rounded corners=2pt, font=\sffamily\scriptsize, minimum width=9.5cm, minimum height=1.6cm, align=left},
  arrow/.style={-{Stealth[scale=0.9]}, thick, color=warningRed}
]

\node[ram] (r) at (0, 2) {
  \textbf{DRAM (Tiny Footprint $\approx 32$ GB)}:\\
  $\bullet$ Compressed Vectors via Product Quantization (PQ) (1-byte codes)\\
  $\bullet$ Vamana Graph Top Entry Points (Inverted Layout)
};

\node[ssd] (s) at (0, -0.4) {
  \textbf{NVMe SSD (Petabyte Scale)}:\\
  $\bullet$ Full Precision 32-bit Floating Point Vectors\\
  $\bullet$ Full Vamana Graph Adjacency List (Single Dispatched Asynchronous IO)
};

\draw[arrow] (r.south) -- (s.north) node[midway, right, font=\tiny\sffamily, color=warningRed] {{1-Hop Asynchronous NVMe Read (\texttt{io\_uring})}};

\end{tikzpicture}
\caption{DiskANN Architecture: Compressed PQ vectors in DRAM guide navigation; only terminal candidates trigger 4KB page reads from NVMe SSD.}
\label{fig:diskann_layout}
\end{figure}

\begin{underthehoodbox}[How DiskANN Achieves Single-Digit Millisecond SSD Search]
DiskANN relies on two innovations:
\begin{enumerate}
  \item \textbf{Vamana Graphs}: Unlike HNSW which uses multiple layers, Vamana is a single-layer flat graph. It introduces $\alpha$-pruned long-range edges ($\alpha \ge 1.2$), ensuring small-world navigation without requiring intermediate layer lookups.
  \item \textbf{Compressed DRAM Guiding}: A 1-byte-per-dimension Product Quantized (PQ) copy of the entire dataset resides in RAM. The beam search navigates the graph in RAM using fast approximate distances; only the actual neighbor node lists are fetched from NVMe via Linux asynchronous I/O (\texttt{io\_uring}).
\end{enumerate}
\end{underthehoodbox}

\section{Engine Internals: Qdrant vs pgvector}

\subsection{Qdrant: Rust-Native Segment Architecture}
Qdrant is written in Rust and designed specifically for vector payloads and dynamic filtering.
\begin{itemize}[leftmargin=18pt]
  \item \textbf{Segmented Storage}: Collections are divided into immutable segments. When a segment reaches capacity, it is frozen and an HNSW graph is built in background worker threads.
  \item \textbf{Filtered HNSW}: Unlike traditional vector databases that perform either pre-filtering (filtering relational metadata first, breaking the graph into disconnected components) or post-filtering (searching top-$k$ vectors and discarding non-matches, causing low recall), Qdrant traverses the HNSW graph while checking a bitset of matching payload IDs \textbf{inside the graph exploration step}.
\end{itemize}

\subsection{pgvector: PostgreSQL Access Methods}
\texttt{pgvector} is a native C extension for PostgreSQL that integrates vector indexing directly into the relational buffer pool and query planner.

\begin{lstlisting}[language=SQL, caption={pgvector Production Configuration and Indexing}]
-- Enable extension
CREATE EXTENSION IF NOT EXISTS vector;

-- Define table with 1536-dimensional unit-normalized embeddings
CREATE TABLE threat_signatures (
    id UUID PRIMARY KEY DEFAULT gen_random_uuid(),
    tenant_id INT NOT NULL,
    observed_at TIMESTAMPTZ NOT NULL,
    embedding vector(1536) NOT NULL
);

-- Production HNSW Index Creation
-- Using cosine distance (<=>) or inner product (<#>)
CREATE INDEX idx_threat_hnsw ON threat_signatures 
USING hnsw (embedding vector_cosine_ops)
WITH (m = 24, ef_construction = 128);

-- Setting runtime search depth
SET hnsw.ef_search = 64;

-- Hybrid Query combining relational filtering and vector search
EXPLAIN ANALYZE
SELECT id, tenant_id, 1 - (embedding <=> '[0.012, -0.045, ...]'::vector) AS similarity
FROM threat_signatures
WHERE tenant_id = 42 AND observed_at >= NOW() - INTERVAL '7 days'
ORDER BY embedding <=> '[0.012, -0.045, ...]'::vector
LIMIT 10;
\end{lstlisting}

\begin{productionwarningbox}[pgvector Shared Buffer Thrashing]
An HNSW index on 1,000,000 vectors with $D=1536$ exceeds 6 GB. If your PostgreSQL \texttt{shared\_buffers} is configured to the default 128 MB or 2 GB, random page reads during HNSW graph traversal will thrash the buffer cache, driving disk I/O to 100\% and spiking query latency from 5ms to 800ms. In production, ensure \texttt{shared\_buffers} or operating system page cache is large enough to pin the entire \texttt{.hnsw} index in memory.
\end{productionwarningbox}

\begin{realworldbox}[Multi-Tenant Vector Isolation Failure]
A healthcare SaaS company implemented a shared vector collection where each patient record had a \texttt{hospital\_id} metadata tag. They queried the top 10 vectors and applied a post-filter: \texttt{WHERE hospital\_id = 'H-901'}.
Because Hospital 901 represented only 0.2\% of the total vectors in the database, the top 10 nearest global vectors never belonged to Hospital 901. The post-filter returned 0 results, creating a severe outage.
Switching to Qdrant's payload-indexed traversal or pgvector's composite index resolved the issue immediately.
\end{realworldbox}

\begin{tryityourselfbox}[Tuning HNSW for Precision vs Speed]
Using pgvector or a local Qdrant instance:
\begin{enumerate}
  \item Ingest 50,000 random vectors of dimension 768.
  \item Compute exact ground-truth top-10 neighbors for 100 test queries using a flat scan.
  \item Benchmark recall@10 and latency (QPS) across $efSearch \in \{8, 16, 32, 64, 128, 256\}$.
  \item Plot the Pareto frontier of Recall vs Query Latency.
\end{enumerate}
\end{tryityourselfbox}

\section{Chapter Summary \& Key Takeaways}
\begin{itemize}[leftmargin=18pt]
  \item \textbf{Normalization Identity}: When vectors are unit-normalized ($\|\vec{v}\|=1$), Cosine distance and Euclidean distance simplify directly to inner products, eliminating square root computations.
  \item \textbf{Curse of Dimensionality}: In spaces with $D \ge 500$, pairwise distance variance vanishes. Traditional spatial trees ($k$-d trees) collapse to linear scans; graph-based heuristics are mandatory.
  \item \textbf{IVF Architecture}: Partitions space into Voronoi cells. It yields low memory overhead but requires careful tuning of $nlist$ and $nprobe$ to prevent boundary recall loss.
  \item \textbf{HNSW Superiority}: Hierarchical skip-list graphs provide logarithmic search time $O(\log N)$ with $>95\%$ recall. $M$, $efConstruction$, and $efSearch$ govern memory, build time, and query speed.
  \item \textbf{DiskANN}: Solves the multi-terabyte RAM bottleneck by keeping compressed Product Quantization codes in DRAM and dispatching single-hop asynchronous I/O reads to NVMe SSDs.
  \item \textbf{Engine Primitives}: Choose Qdrant when workloads demand complex payload filtering and high-concurrency ingestion; choose pgvector when vector data must live transactionally alongside relational ACID tables.
\end{itemize}

\section{Exercises}

\subsection*{Mechanics}
\textbf{Exercise 3.1}: Prove that for any two vectors $\vec{u}, \vec{v} \in \mathbb{R}^D$ where $\|\vec{u}\|_2 = 1$ and $\|\vec{v}\|_2 = 1$, the Euclidean distance squared is strictly given by $\|\vec{u} - \vec{v}\|_2^2 = 2 - 2\langle \vec{u}, \vec{v} \rangle$.

\textbf{Exercise 3.2}: Calculate the memory footprint (in gigabytes) required to store an HNSW index of $10,000,000$ vectors with dimension $D=1536$ (stored as 32-bit floats), assuming average node degree $M=32$ where each edge is stored as a 64-bit integer pointer.

\subsection*{Debugging}
\textbf{Exercise 3.3}: An engineer complains that their \texttt{pgvector} HNSW queries take 350ms despite having created an HNSW index. Upon inspecting \texttt{EXPLAIN ANALYZE}, they see:
\texttt{Seq Scan on embeddings (cost=0.00..84500.00 rows=10 width=6144)}. Explain why the query planner chose a sequential scan over the HNSW index.

\subsection*{Architecture}
\textbf{Exercise 3.4}: Design a vector search infrastructure for an e-commerce catalog of 500 million product items. The system must support real-time price updates (1,000 updates/sec), stock-availability filtering, and sub-20ms vector search latency under a budget that prohibits provisioning 2 Terabytes of RAM.

\section{Detailed Solutions}

\subsection*{Solution 3.1}
By definition of the $L_2$ norm:
\begin{align*}
\|\vec{u} - \vec{v}\|_2^2 &= \sum_{i=1}^D (u_i - v_i)^2 = \sum_{i=1}^D (u_i^2 - 2u_i v_i + v_i^2) \\
&= \sum_{i=1}^D u_i^2 + \sum_{i=1}^D v_i^2 - 2 \sum_{i=1}^D u_i v_i \\
&= \|\vec{u}\|_2^2 + \|\vec{v}\|_2^2 - 2\langle \vec{u}, \vec{v} \rangle
\end{align*}
Substituting the unit norm condition $\|\vec{u}\|_2 = \|\vec{v}\|_2 = 1$:
\[
\|\vec{u} - \vec{v}\|_2^2 = 1 + 1 - 2\langle \vec{u}, \vec{v} \rangle = 2 - 2\langle \vec{u}, \vec{v} \rangle \quad \blacksquare
\]

\subsection*{Solution 3.2}
For each vector:
\begin{enumerate}
  \item Vector Float32 storage: $1536 \times 4 \text{ bytes} = 6,144 \text{ bytes}$.
  \item HNSW edge pointers: Across all layers, average degree is approximately $2M$ on layer 0 and $M$ on upper layers. Approximating with $M = 32$ bidirectional links on layer 0:
  $32 \text{ edges} \times 8 \text{ bytes (int64)} = 256 \text{ bytes}$.
  Upper layer overhead adds roughly 10\%: $256 \times 1.1 \approx 282 \text{ bytes}$.
  \item Total per vector $\approx 6,144 + 282 = 6,426 \text{ bytes}$.
\end{enumerate}
For $10,000,000$ vectors:
\[
\text{Total Memory} = 10^7 \times 6,426 \text{ bytes} = 64,260,000,000 \text{ bytes} \approx 64.26 \text{ GB}
\]
(If product quantization is not applied, 64.3 GB of RAM is required strictly for index memory).

\subsection*{Solution 3.3}
The PostgreSQL query planner bypassed the HNSW index due to one of three common causes:
\begin{enumerate}
  \item \textbf{Distance operator mismatch}: The index was created with \texttt{vector\_cosine\_ops} (requiring \texttt{\textless{}=\textgreater{}}), but the query used Euclidean distance \texttt{\textless{}-\textgreater{}} or inner product \texttt{\textless{}\#\textgreater{}}.
  \item \textbf{Missing ORDER BY / LIMIT}: HNSW indices are only utilized when an explicit \texttt{ORDER BY embedding \textless{}=\textgreater{} query\_vector LIMIT k} is present.
  \item \textbf{Low Planner Cost Setting}: The table statistics are outdated, or \texttt{enable\_seqscan} is on while the table contains too few rows (under 10,000) where the planner estimates a sequential scan is faster than random page lookups.
\end{enumerate}

\subsection*{Solution 3.4}
Architecture:
\begin{enumerate}
  \item \textbf{DiskANN / Vamana Index on NVMe}: Store full-resolution vectors on high-IOPS NVMe SSDs. Retain 1-byte Product Quantized (PQ) compressed representations in DRAM (consuming only $500\text{M} \times 64 \text{ bytes} \approx 32 \text{ GB}$ of RAM).
  \item \textbf{Decoupled Relational / In-Memory Filter}: Fast-changing attributes (price, stock status) are kept in Redis bitsets or memory-mapped relational tables, decoupled from the vector index.
  \item \textbf{Two-Stage Filter Execution}: Beam search traverses the SSD Vamana graph. At each hop, the candidate node ID is intersected with the in-memory stock bitset before expanding.
\end{enumerate}

\section{Further Reading \& Citations}
\begin{itemize}[leftmargin=18pt]
  \item Malkov, Y. A., \& Yashunin, D. A. (2018). \emph{Efficient and Robust Approximate Nearest Neighbor Search Using Hierarchical Navigable Small World Graphs}. IEEE Transactions on Pattern Analysis and Machine Intelligence, 42(4), 824-836.
  \item Subramanya, S. J., et al. (2019). \emph{DiskANN: Fast Accurate Billion-point Nearest Neighbor Search on a Single Node}. Advances in Neural Information Processing Systems (NeurIPS).
  \item Jégou, H., Douze, M., \& Schmid, C. (2010). \emph{Product Quantization for Nearest Neighbor Search}. IEEE Transactions on Pattern Analysis and Machine Intelligence, 33(1), 117-128.
\end{itemize}
